\documentclass[11pt]{article}

\usepackage[margin=1in]{geometry}
\usepackage{amsmath,amssymb,mathtools}
\usepackage{microtype}
\usepackage[hidelinks]{hyperref}

\newcommand{\bmat}[1]{\begin{bmatrix}#1\end{bmatrix}}
\newcommand{\ip}[2]{\left\langle #1,#2\right\rangle}

\begin{document}
\begin{center}
 {\LARGE Triangularization and Duality of the Normalized IQC System}
\end{center}
\vspace{1em}

Let $P$ be the primal input--output operator and let
\begin{equation*}
 \Psi=\bmat{\Psi_{11}&\Psi_{12}\\\Psi_{21}&\Psi_{22}},
 \qquad
 J_{z,w}=\bmat{I_z&0\\0&-I_w},
 \qquad
 J_{w,z}=\bmat{I_w&0\\0&-I_z}.
\end{equation*}
Assume that $\Psi$ is boundedly and causally invertible.  To use the
triangularization of Yuan and Wu, assume additionally that $\Psi_{22}$ has a
bounded causal inverse.  Finally, assume that
\begin{equation*}
 F_0:=\Psi_{21}P+\Psi_{22}
\end{equation*}
has a bounded causal inverse.  This is the input--output counterpart of the
current assumption in Lemma~1,
$(\mathcal D_\Theta^{wz}\mathcal D+\mathcal D_\Theta^w)^{-1}\in\Pi_4$.
Throughout, $\top$ denotes the causal dual of an input--output operator;
$*$ is retained for realization operators.

\section*{Triangularize the factor and normalize the primal graph}

Following Yuan and Wu, introduce the invertible signal transformation
\begin{equation*}
 L_\Psi:=\bmat{I&0\\\Psi_{21}&\Psi_{22}},
 \qquad
 \bmat{z\\\tilde w}=L_\Psi\bmat{z\\w}.
\end{equation*}
Moving this transformation from the factor to the interconnection gives
\begin{equation*}
 \Psi_\triangle:=\Psi L_\Psi^{-1}
 =\bmat{
 \Psi_{11}-\Psi_{12}\Psi_{22}^{-1}\Psi_{21}
 &\Psi_{12}\Psi_{22}^{-1}\\
 0&I
 }.
\end{equation*}
Hence the triangular form is obtained by an invertible change of internal
signals; it is not an additional restriction on the original factor.
Moreover,
\begin{equation*}
 \Psi\bmat{z\\w}
 =\Psi_\triangle\bmat{z\\\tilde w}.
\end{equation*}

On the graph $z=Pw$, the same transformation gives
\begin{equation*}
 L_\Psi\bmat{P\\I}w
 =\bmat{P\\F_0}w
 =\bmat{PF_0^{-1}\\I}F_0w.
\end{equation*}
Consequently,
\begin{equation*}
 \Psi\bmat{P\\I}w
 =\Psi_\triangle\bmat{PF_0^{-1}\\I}F_0w
 =\bmat{G_0\\I}F_0w,
\end{equation*}
where
\begin{equation*}
 G_0=(\Psi_{11}P+\Psi_{12})
     (\Psi_{21}P+\Psi_{22})^{-1}.
\end{equation*}
Since $\tilde w=F_0w$, the normalized system is
\setcounter{equation}{20}
\begin{equation}
 \widetilde z=G_0\widetilde w.
 \label{eq:normalized-system}
\end{equation}

The primal IQC is now an ordinary gain condition:
\begin{align*}
 \mathcal I_T(w)
 &:=
 \ip{p_T\Psi\bmat{P\\I}w}
 {J_{z,w}p_T\Psi\bmat{P\\I}w}_{L_\mathbf H}\\
 &=\|p_TG_0F_0w\|^2
   -\|p_TF_0w\|^2.
\end{align*}
If $\|G_0\|\leq\rho<1$, the bounded causal invertibility of $F_0$ assumed at
the start gives
\begin{equation*}
 \mathcal I_T(w)
 \leq-\frac{1-\rho^2}{\|F_0^{-1}\|^2}
       \|p_Tw\|^2.
\end{equation*}
In a state-space representation, $L_\Psi$ induces the corresponding
invertible congruence transformation between the original filtered
realization and the normalized realization.

\section*{Obtain the paired dual transformation}

Define
\begin{equation*}
 \widehat J:=\bmat{0&I\\-I&0},
 \qquad
 \mathbf D(\Psi):=(\widehat J^*\Psi^\top\widehat J)^{-1}.
\end{equation*}
Calculate the dual filter directly.  Since
\begin{equation*}
 \widehat J^*\Psi^\top\widehat J
 =\bmat{
 \Psi_{22}^\top&-\Psi_{12}^\top\\
 -\Psi_{21}^\top&\Psi_{11}^\top
 },
\end{equation*}
and
\begin{equation*}
 (\Psi_{\triangle,11})^\top
 =\Psi_{11}^\top
  -\Psi_{21}^\top(\Psi_{22}^\top)^{-1}\Psi_{12}^\top,
\end{equation*}
the block inverse formula gives
\begin{equation*}
 \mathbf D(\Psi)
 =\bmat{
 (\Psi_{22}^\top)^{-1}
 +(\Psi_{22}^\top)^{-1}\Psi_{12}^\top
  ((\Psi_{\triangle,11})^\top)^{-1}
  \Psi_{21}^\top(\Psi_{22}^\top)^{-1}
 &
 (\Psi_{22}^\top)^{-1}\Psi_{12}^\top
  ((\Psi_{\triangle,11})^\top)^{-1}
 \\
 ((\Psi_{\triangle,11})^\top)^{-1}
  \Psi_{21}^\top(\Psi_{22}^\top)^{-1}
 &
 ((\Psi_{\triangle,11})^\top)^{-1}
 }.
\end{equation*}
The desired paired factorization can now be read off directly:
\begin{equation*}
 \underline\Psi_\triangle
 :=\bmat{
 (\Psi_{22}^\top)^{-1}
 &(\Psi_{22}^\top)^{-1}\Psi_{12}^\top\\
 0&I
 },
 \qquad
 \underline L_\Psi
 :=\bmat{
 I&0\\
 \mathbf{D}(\Psi)_{21}
 &\mathbf{D}(\Psi)_{22}
 }.
\end{equation*}

\section*{Normalize the dual graph}

Write $S:=\Psi_{\triangle,11}$.  The lower block of $\underline L_\Psi$
defines the dual normalizer
\begin{equation*}
 \underline\Gamma_0
 :=\mathbf D(\Psi)_{21}P^\top+\mathbf D(\Psi)_{22}
 =S^{-\top}\left(I+\Psi_{21}^\top
  \Psi_{22}^{-\top}P^\top\right).
\end{equation*}
Equivalently,
\begin{equation*}
 \underline\Gamma_0^\top
 =\left(I+P\Psi_{22}^{-1}\Psi_{21}\right)S^{-1}.
\end{equation*}
Because $F_0=\Psi_{22}(I+\Psi_{22}^{-1}\Psi_{21}P)$ is invertible, the
push-through rule shows that $I+P\Psi_{22}^{-1}\Psi_{21}$ is invertible.
Hence $\underline\Gamma_0$ is invertible, with
\begin{equation*}
 \underline\Gamma_0^{-\top}
 =S\left(I+P\Psi_{22}^{-1}\Psi_{21}\right)^{-1}.
\end{equation*}

On the dual graph $\underline z=P^\top\underline w$, the transformation
$\underline L_\Psi$ now gives
\begin{equation*}
 \underline L_\Psi\bmat{P^\top\\I}\underline w
 =\bmat{P^\top\\\underline\Gamma_0}\underline w
 =\bmat{P^\top\underline\Gamma_0^{-1}\\I}
  \underline\Gamma_0\underline w.
\end{equation*}
Consequently,
\begin{align*}
 \mathbf D(\Psi)\bmat{P^\top\\I}\underline w
 &=\underline\Psi_\triangle
   \bmat{P^\top\underline\Gamma_0^{-1}\\I}
   \underline\Gamma_0\underline w\\
 &=\bmat{\underline G_0\\I}
   \underline\Gamma_0\underline w,
\end{align*}
where
\begin{equation*}
 \underline G_0
 =(\Psi_{22}^\top)^{-1}
  \left(P^\top\underline\Gamma_0^{-1}+\Psi_{12}^\top\right).
\end{equation*}

The same push-through identity gives
\begin{equation*}
 (I+P\Psi_{22}^{-1}\Psi_{21})^{-1}P\Psi_{22}^{-1}
 =P(\Psi_{21}P+\Psi_{22})^{-1}=PF_0^{-1}.
\end{equation*}
It therefore identifies the normalized dual operator in three lines:
\begin{align*}
 \underline G_0^\top
 &=\underline\Gamma_0^{-\top}P\Psi_{22}^{-1}
   +\Psi_{12}\Psi_{22}^{-1}\\
 &=\left(SP
   +\Psi_{12}\Psi_{22}^{-1}F_0\right)F_0^{-1}\\
 &=(\Psi_{11}P+\Psi_{12})F_0^{-1}
 =G_0.
\end{align*}
Hence, with $\tilde{\underline w}:=\underline\Gamma_0\underline w$, the
normalized dual system is
\begin{equation*}
 \tilde{\underline z}=G_0^\top\tilde{\underline w}.
\end{equation*}
Equivalently,
\begin{equation*}
 G_0^\top
 =(P^\top\Psi_{21}^\top+\Psi_{22}^\top)^{-1}
  (P^\top\Psi_{11}^\top+\Psi_{12}^\top),
\end{equation*}
where the first factor is $(F_0^\top)^{-1}=(F_0^{-1})^\top$.

The dual IQC is now the same ordinary gain condition:
\begin{align*}
 \underline{\mathcal I}_T(\underline w)
 &:={}
 \ip{p_T\mathbf D(\Psi)\bmat{P^\top\\I}\underline w}
 {J_{w,z}p_T\mathbf D(\Psi)\bmat{P^\top\\I}\underline w}_{L_\mathbf H}\\
 &=\|p_TG_0^\top\underline\Gamma_0\underline w\|^2
   -\|p_T\underline\Gamma_0\underline w\|^2.
\end{align*}
Since $\|G_0^\top\|=\|G_0\|\leq\rho<1$, bounded causal invertibility of
$\underline\Gamma_0$ gives
\begin{equation*}
 \underline{\mathcal I}_T(\underline w)
 \leq-\frac{1-\rho^2}{\|\underline\Gamma_0^{-1}\|^2}
       \|p_T\underline w\|^2.
\end{equation*}
For the sign comparison, no strict upper bound on $\rho$ is needed.  For any
common gain bound $\|G_0\|=\|G_0^\top\|\leq\rho$, define the pullback
coefficients
\begin{equation*}
 \varepsilon_{\mathrm p}:=
 \frac{1-\rho^2}{\|F_0^{-1}\|^2},
 \qquad
 \varepsilon_{\mathrm d}:=
 \frac{1-\rho^2}{\|\underline\Gamma_0^{-1}\|^2}.
\end{equation*}
Because both normalizers are boundedly invertible, their inverse norms are
finite and strictly positive.  Consequently,
\begin{equation*}
 \operatorname{sgn}(\varepsilon_{\mathrm p})
 =\operatorname{sgn}(1-\rho^2)
 =\operatorname{sgn}(\varepsilon_{\mathrm d}).
\end{equation*}
In particular, $\varepsilon_{\mathrm p}>0$ if and only if
$\varepsilon_{\mathrm d}>0$.  A strict bound $\rho<1$ is needed only to
conclude that this common sign is positive, in which case the coefficients are
valid strict IQC margins.  Thus primal and dual strictness are equivalent,
although the two positive margins need not have the same magnitude.

\end{document}
